\section{The \emph{Celestial Hierarchy}}
\label{sec:celestial-hierarchy}

The author received as Dionysius the Areopagite speaks under the identity of Paul's Athenian convert (Acts 17:34). Aquinas receives him within that apostolic tradition. Historical criticism places the corpus in the late fifth or early sixth century; its engagement with Proclus is an important reason for distinguishing its author from the first-century convert.\footnote{Eric Perl, ``Pseudo-Dionysius the Areopagite,'' \emph{The Cambridge History of Philosophy in Late Antiquity}, ch.~42, published online 28 May 2011, publisher's summary, \url{https://doi.org/10.1017/CHOL9780521194846.014}; Benedict XVI, general audience, 14 May 2008, on the pseudonym, sixth-century setting, and relation to Proclus.}

\subsection{I: illumination that gathers its recipients}
\emph{That every divine illumination, whilst going forth lovingly to the objects of its forethought under various forms, remains simplex. Nor is this all. It also unifies the things illuminated.}\footnote{CH I.1--3; John Parker (trans.), \emph{On the Heavenly Hierarchy}, in \emph{The Works of Dionysius the Areopagite}, part~II (London and Oxford: James Parker and Co., 1899), pp.~1--4; PG~3, cols.~119--124.}

The opening movement is from the Father of lights through Jesus, the light who gives access to the Father. Every ``gift of goodness'' (I.1, p.~1) comes from that source. Multiplicity belongs to the gifts and their recipients; God does not break into pieces as creatures participate in his light. Illumination draws those recipients into unity with their source and with one another. The book therefore begins with divine generosity, before it distinguishes the ranks that receive it.

Human beings need sensible mediation. Scriptural images and sacred rites accommodate an embodied intelligence, drawing it through visible arrangements toward intelligible truth. They neither imprison God in an image nor become dispensable merely because the intended truth is spiritual. The heavenly hierarchy supplies a pattern for this ascent; the earthly hierarchy expresses that pattern under conditions appropriate to us.

Aquinas expressly uses this chapter in \ST{I q.~111 a.~1} to explain why an angel illuminating a human being presents truth under sensible likenesses. His account in \ST{I q.~1 a.~9} similarly makes metaphor suitable to human cognition. Reception here sharpens the difference between recipients: an angel does not need bodily pictures in order to understand, while a human learner naturally begins from sense. Divine accommodation is not a defect in revelation. It is one form of the goodness with which the chapter opens.

\subsection{II: why the symbols can be strange}
\emph{That Divine and Heavenly things are appropriately revealed, even through dissimilar symbols.}\footnote{CH II.1--5; Parker, pp.~4--13; PG~3, cols.~135--146.}

Dionysius confronts the difficulty that sacred representations can look unworthy of heavenly intelligence. Scripture presents beings ``many-footed and many-faced'' (II.1, p.~4), animals, fire, and material implements. Why not reserve beautiful human forms for beautiful spiritual realities? Because a dignified image can tempt its reader to stop at the image. An obviously incongruous likeness announces its own inadequacy and prompts the mind to seek what it signifies.

The method does not make ugliness divine or render interpretation arbitrary. A creature can supply a likeness because it receives some good from its cause; its limitations must then be denied of the transcendent reality. Thus a bodily appetite may signify spiritual desire without making an angel hungry. Affirmation and negation work together: one identifies a relevant perfection, the other removes its bodily mode. The same image may signify different things when its relations and context differ.

In \ST{I q.~1 a.~9 ad~3}, Aquinas repeats this preference for less noble figures as protection against literal misunderstanding.\footnote{The received English \emph{Summa} reference labels its Dionysian source chapter~I; the developed argument occurs in CH II.} Aquinas also explains that bodily organs ascribed to angels signify spiritual powers (\ST{I q.~51 a.~3 ad~2}). Both discussions restrain the imagination without treating the scriptural image as worthless.

\subsection{III: hierarchy as participation and service}
\emph{What is Hierarchy? and what the use of Hierarchy?}\footnote{CH III.1--3; Parker, pp.~13--16; PG~3, cols.~163--168.}

Hierarchy is ``a sacred order and science and operation'' (III.1, p.~13). The definition joins an arrangement of persons, a reception of truth, and an activity. Its end is likeness to God and union with him as far as the recipient can attain it. Rank has its meaning inside that end. A superior becomes an image of divine generosity by communicating what has been received; the subordinate becomes capable of receiving and communicating in turn.

Purification, illumination, and perfection describe the hierarchy's activity. Purification removes impediments to reception; illumination imparts light; perfection establishes the recipient in the knowledge and activity proper to its measure. God is their source without undergoing them. The same language therefore changes its force when used of God, of the heavenly spirits, and of human beings being freed from sin. The definition cannot simply be transferred unchanged to every subject.

Aquinas makes that boundary explicit in \ST{I q.~108 a.~1, body}: the divine Persons do not form a hierarchy, since no Person is purified, enlightened, or perfected by another. Among holy angels, purification can mean removal of ignorance concerning something newly revealed, without the cleansing of moral fault (\ST{I q.~106 a.~2 ad~1}). These distinctions preserve the productive insight of Dionysius's triad while excluding an inequality in the Trinity or a residual stain in the blessed.

\subsection{IV: why the heavenly spirits are messengers}
\emph{What is meant by the appellation ``Angels?''}\footnote{CH IV.1--4; Parker, pp.~16--21; PG~3, cols.~177--182.}

All creatures exist ``through goodness'' (IV.1, p.~17). Dionysius distinguishes participation in being, in life, and in wisdom, then places the intellectual heavenly orders among the creatures most receptive to divine likeness. Their proximity is a mode of participation, not a position alongside God as another uncreated principle. Their name follows their work: they announce and manifest divine illumination to those below them.

This explains the angelic mediation of scriptural revelations. A manifestation can truly disclose God while its created messenger and sensible form remain distinct from God's essence. The chapter moves from the Law and the prophets to the announcements surrounding the Incarnation. The incarnate Jesus accepts angelic ministry within his human life even while he remains the source of their order. Mediation serves his saving action; it does not place him beneath a heavenly administration.

Aquinas cites chapter~IV when affirming angelic illumination of human beings (\ST{I q.~111 a.~1, sed contra}) and the orderly communication of heavenly knowledge (\ST{I q.~106 a.~3, sed contra}). But this does not require a superior angel to stand between every blessed intellect and God as an object of vision. In \ST{I q.~106 a.~1 ad~1}, all the blessed see the divine essence immediately; higher angels illuminate lower ones about the divine works and mysteries. Messenger language must therefore be applied to the kind of knowledge actually in question.

\subsection{V: a common name and a particular order}
\emph{For what reason all the Heavenly Beings are called, in common, Angels.}\footnote{CH V, unnumbered; Parker, pp.~21--22; PG~3, cols.~195--196.}

``Angel'' names both the whole heavenly company and its lowest order. Dionysius resolves the apparent ambiguity through the relation between superior and subordinate capacities: ``superior ranks possess the illuminations and powers of their subordinates'' (p.~22). The higher spirits can announce God's gifts; the lowest order does not therefore possess every excellence signified by the names seraphim or thrones. A shared activity need not eliminate a distinctive eminence.

The argument also makes the lowest order intelligible. It completes the ordered communication toward humanity, receiving what it passes on. Its name emphasizes the ministry most immediately apparent to us. No inference follows that other angels are inactive or that the lowest are outside divine communion. Their work belongs to the same movement of reception and communication described in chapter~III.

Aquinas offers a closely related account in \ST{I q.~108 a.~5}: a perfection can belong to a subject properly, by participation, or in a more excellent manner. Thus the order named angels possesses the common ministerial perfection under the ordinary mode designated by its name, while higher orders possess it together with a more eminent character. The distinction is useful whenever a biblical passage says simply ``angel'': the word alone does not establish which of the nine orders is intended.

\subsection{VI: nine names and the limits of the map}
\emph{Which is the first Order of the Heavenly Beings? which the middle? and which the last?}\footnote{CH VI; Parker, pp.~23--24; PG~3, cols.~199--202, divided there into sections~1--2.}

God ``alone distinctly knows'' the heavenly arrangements in their exact character (VI.1, p.~23). Dionysius allows their heavenly recipients a knowledge of their own gifts, while restricting his account to what sacred teaching communicates to human beings. This limitation precedes the famous scheme. The nine names are arranged in three triads: seraphim, cherubim, thrones; dominions, virtues, powers; principalities, archangels, angels. Parker's different English words for the middle triad require special care in chapter~VIII.

The chapter's brevity should not make its opening reserve disappear. Nine orders are not nine individual angels, and a human account of their common characters is not a complete description of every spirit. Nor does the presence of nine names scattered through Scripture itself give the reader their full sequence. Dionysius is presenting an interpretation of their arrangement, whose rationale the following chapters develop.

Aquinas expressly receives the threefold ordering in \ST{I q.~108 aa.~1--2}. He then observes that we know the orders imperfectly and distinguish them through general offices; God knows distinctions hidden from us, and each angel has its own particular rank (\ST{I q.~108 a.~3}). The scholastic account therefore need not turn the diagram into an exhaustive census of heavenly difference. Its general order and its admitted limits belong together.

\subsection{VII: fire, knowledge, and the bearing of God}
\emph{Concerning the Seraphim and Cherubim and Thrones, and concerning their first Hierarchy.}\footnote{CH VII.1--4; Parker, pp.~24--31; PG~3, cols.~205--212.}

Seraphim signify ``kindling or burning'' (VII.1, p.~24): an unceasing movement toward God, fervor communicated to others, and a power of purification. Cherubim signify abundant knowledge and wisdom, receptive contemplation, and the generous imparting of what is known. Thrones signify stability, freedom from debasing attachment, and readiness to receive God. Their bearing is intellectual and spiritual; the imagery does not make God a body needing a seat.

The first hierarchy receives its initiation from God with a proximity proper to it. Yet it remains a recipient. Dionysius explains its purification without imagining that it first needed release from sensual sin. Greater disclosure removes ignorance, while illumination and perfection enlarge the creature's participation. The questions attributed to heavenly beings and their hymns disclose both the reception of mysteries and their return to God in praise.

Aquinas draws directly on these names in \ST{I q.~108 a.~5 ad~5--6}. Love's fervor, fullness of knowledge, and receptivity to judgment express distinguishable excellences without excluding charity, knowledge, or obedience from the other orders. He also cites this chapter in explaining angelic illumination and immediately distinguishes the vision of God's essence from knowledge of God's works (\ST{I q.~106 a.~1}). Finally, the questions addressed by lower to higher angels establish that speech is broader than illumination: a question can move upward without making its speaker the other's teacher (\ST{I q.~107 a.~2}).

\subsection{VIII: authority freed from servility and tyranny}
\emph{Concerning Lordships and Powers and Authorities, and concerning their middle Hierarchy.}\footnote{CH VIII.1--2; Parker, pp.~31--35; PG~3, cols.~237--242.}

Here Parker's Lordships correspond to dominions, Powers to virtues, and Authorities to powers in the usual Latin-derived English list. Confusing the two uses of ``powers'' exchanges different orders. Dionysius describes lordship through ``unslavish elevation'' (VIII.1, p.~31): freedom from servile deformation and a likeness to divine sovereignty. It is an ordered liberty from what degrades, rather than a license to exercise arbitrary will.

The virtues express vigorous, unwearied readiness for divine action. Their strength does not draw them away from their giver; it makes them receptive and capable of strengthening others. The powers express a well-ordered exercise of authority that brings others toward God without tyranny. The common character of this triad is a received capacity for governance. Its members do not originate the divine intention they administer.

Dionysius supports mediated illumination with prophetic visions in which heavenly messengers communicate with one another. Aquinas cites this chapter for the claim that one angel enlightens another (\ST{I q.~106 a.~1, sed contra}). In explaining the names, he relates virtues to strength and powers to the ordering of action (\ST{I q.~108 a.~5 ad~1--3}). His later distinction between commanding and executing a ministry prevents dominion from being confused with personally performing every external mission (\ST{I q.~112 a.~4}). The activity of a superior can be real precisely through ordered cooperation.

\subsection{IX: the last triad and providence for the nations}
\emph{Concerning the Principalities, Archangels, and Angels, and concerning their last Hierarchy.}\footnote{CH IX.1--4; Parker, pp.~35--40; PG~3, cols.~257--262.}

Principalities have a ``princely and leading function'' (IX.1, p.~36), turned first toward God and then exercised toward others. Archangels occupy a middle position, sharing the directing character of principalities and the announcing ministry of angels. Angels complete the heavenly ranks and stand nearest the human order in their ministry. The nearest is not therefore the most divine; nearness here concerns the communication reaching us.

Dionysius extends the discussion to the guardianship of nations. Israel's relation to God does not mean that other peoples have different supreme creators or that the one providence is divided into rival territories. The same God governs all. The failure of a nation to receive illumination is not proof of a failure of goodness in its appointed angel: freedom can resist the gift. The examples of Melchizedek and revelations involving Egypt and Babylon prevent the account from becoming a claim that divine care stops at one people's border.

Aquinas uses chapter~IX in explaining principalities and the intermediate office of archangels (\ST{I q.~108 a.~5 ad~3--4}). His distinction between the guardianship of a multitude and of an individual develops the implications of corporate ministry (\ST{I q.~113 a.~3}). The limit of such ministry follows from a separate determination: God alone moves the will inwardly as its cause, while an angel can persuade through an apprehended good (\ST{I q.~111 a.~2}). Providence through a guardian does not abolish the guarded person's agency.

\subsection{X: distinction within the distinctions}
\emph{A Repetition and Summary of the Angelic discipline.}\footnote{CH X.1--3; Parker, pp.~41--42; PG~3, cols.~271--274.}

The first hierarchy's illuminations are ``more hidden and more manifest'' (X.1, p.~41). The expressions are compatible because they answer different questions. The illumination is less accessible to a bodily imagination and less divided into forms; it is also more immediate and unveiled to its spiritual recipients. To call it hidden is not to say that the first hierarchy understands less than the last.

The chapter summarizes the movement through the three hierarchies, then prevents the summary from becoming rigid. First, middle, and last can also be distinguished within an order and within the powers of an individual spirit. Receiving and communicating are not exhausted by assigning one name to each of nine boxes. Every creature has a determinate capacity and relation to what perfects it; only God is the source of perfection without receiving it.

Aquinas cites chapter~X in \ST{I q.~108 a.~3 ad~1} when acknowledging higher, middle, and lower members within the same order. Common office does not require equality in every respect. His account of the communication of illumination supplies a related safeguard: a superior communicates the truth it knows according to the lower recipient's capacity, without thereby making their modes of understanding identical (\ST{I q.~106 a.~4}). The unity of the hierarchy is a shared turning to God through differentiated gifts, not the disappearance of the persons and capacities that receive them.

\subsection{XI: the common power of an intellectual creature}
\emph{For what reason all the Heavenly Beings, in common, are called Heavenly Powers.}\footnote{CH XI.1--2; Parker, pp.~42--44; PG~3, cols.~283--286.}

The explanation of the common name ``angels'' cannot simply be repeated for ``powers.'' Powers are not the lowest order: principalities, archangels, and angels stand below the order that Parker calls Powers, or virtues. If a higher rank contains a lower excellence, that principle alone cannot explain why a lower spirit also receives this higher name. Dionysius therefore distinguishes the specific name from the general character of an intellectual creature.

The relevant analysis is into ``essence, and power, and energy'' (XI.2, p.~43). Every heavenly intellect is something, possesses a capacity to act, and exercises an activity. It may consequently be called a power in a general sense without possessing the distinctive excellence of the order named Powers. The same word functions at different levels of predication. Preserving those levels avoids both an artificial multiplication of orders and an erasure of their differences.

Aquinas expressly cites chapter~XI in \ST{I q.~54 a.~3, sed contra}, where the question is whether an angel's intellectual power is identical with its essence. He answers that created intellectual capacity must be distinguished from the substance that has it. He also uses the chapter to explain the generic and special senses of virtue or power (\ST{I q.~108 a.~5 ad~1}). A problem of heavenly vocabulary thus contributes to a precise metaphysical boundary: the simplicity proper to God cannot be inferred merely from an angel's immateriality.

\subsection{XII: when a human being is called an angel}
\emph{Why the Hierarchs amongst men are called Angels.}\footnote{CH XII.1--3; Parker, pp.~44--45; PG~3, cols.~291--294.}

The scriptural application of messenger language to a human religious minister does not dissolve the difference between human and angelic nature. Dionysius explains it by participation in an announcing and sanctifying office. A recipient may bear a name derived from a superior perfection while possessing that perfection in a limited mode. The same principle governs the scriptural naming of holy persons as gods: participation never makes its recipient the uncreated source.

The orders are joined in ``one harmonious and binding fellowship'' (XII.2, p.~44). This fellowship is more than a shared vocabulary. The lower receives real gifts present in a higher manner above it. Yet the mode of possession matters. To acknowledge a participation in wisdom is not to claim equality in the comprehensiveness with which wisdom is possessed. Communion and difference are mutually intelligible within this account.

Aquinas cites this chapter in \ST{I q.~55 a.~3, sed contra} when distinguishing the more universal intelligible forms of higher angels from the more divided knowledge of lower ones. He also considers it in \ST{I q.~106 a.~4}: the distinction concerns the mode of receiving illumination, not a refusal by the superior to communicate its light. A further comparison concerns humanity's final destiny. Human beings can be associated with angelic orders through gifts of grace without becoming angels by nature (\ST{I q.~108 a.~8}). Neither ecclesial office nor heavenly fellowship requires a change of species.

\subsection{XIII: who purified Isaiah?}
\emph{For what reason the Prophet Isaiah is said to have been purified by the Seraphim.}\footnote{CH XIII.1--4; Parker, pp.~46--53; PG~3, cols.~299--308.}

Isaiah's vision presses a difficulty against ordered mediation. A seraph appears to perform an immediate ministry to a human prophet, although the seraphim belong to the first hierarchy. Dionysius offers two explanations. The minister could be a lower angel called a seraph because the purification is accomplished ``through fire'' (XIII.2, p.~46). Alternatively, the lower minister could attribute the action to the superior seraphic power through which it has received the ability to purify.

The second explanation receives an extended account. As illumination reaches different recipients according to their capacities, the effect can be attributed to the first created bearer of a communicated perfection while being executed by a subordinate. Above both stands God, the source of purification itself. The account seeks to preserve the truth of the prophetic language and the dependence of secondary agents, rather than assign independent purifying powers to competing beings. Dionysius ends by welcoming a better explanation if one can be given.

Aquinas receives both resolutions in \ST{I q.~112 a.~2 ad~2}, within his judgment that the higher orders do not personally undertake exterior missions. Gregory's reserve at \emph{Homilies on the Gospels} XXXIV.13 shows that this conclusion was not the only patristic handling of the evidence. Dionysius's alternatives preserve the seraphic character of the purification through either the minister's office or the higher power it communicates. Aquinas organizes both within his account of ordered mission.

\subsection{XIV: a multitude beyond our reckoning}
\emph{What the traditional number of the Angels signifies.}\footnote{CH XIV, unnumbered; Parker, pp.~53--54; PG~3, cols.~321--322.}

The thousands and ten thousands of the heavenly court express a multitude that ``cannot be numbered by us'' (p.~53). Dionysius does not calculate a total by multiplying the terms of Daniel's vision. The images indicate the greatness of the hosts and the inadequacy of our reckoning. Their number and ordered distinctions remain comprehended by their Creator, whose knowledge is not limited by our ability to count.

The qualification ``by us'' matters. An innumerable multitude in this sense need not be an actually infinite collection. Nor does the passage license an exact ratio between angelic orders, an assignment of a fixed number to each nation, or a claim that only those explicitly counted belong to the heavenly court. Number here serves the revelation of abundance under divine order.

Aquinas cites this chapter in \ST{I q.~50 a.~3}, concluding that the angels are exceedingly numerous and exceed corporeal substances in number. His accompanying argument from the perfection of the universe and the excellence of immaterial creatures is philosophical reasoning, not another numerical datum supplied by Daniel. In \ST{I q.~112 a.~4 ad~2}, he also considers the different ways Gregory and Dionysius read the numbers of those who minister and those who stand before God. The textual abundance can be affirmed without pretending that either treatment provides a census accessible to us.

\subsection{XV: reading the forms without imprisoning the spirits}
The final chapter's printed title supplies its own inventory.\footnote{CH XV.1--9; Parker, pp.~54--66; PG~3, cols.~325--340. Parker's title reads: ``What are the morphic likenesses of the Angelic Powers? what the fiery? what the anthromorphic? what are the eyes? what the nostrils? what the ears? what the mouths? what the touch? what the eyelids? what the eyebrows? what the prime? what the teeth? what the shoulders? what the elbows and the hands? what the heart? what the breasts? what the back? what the feet? what the wings? what the nakedness? what the robe? what the shining raiment? what the sacerdotal? what the girdles? what the rods? what the spears? what the battle-axes? what the measuring lines? what the winds? what the clouds? what the brass? what the electron? what the choirs? what the clapping of hands? what the colours of different stones? what the appearance of the lion? what the appearance of the ox? what the appearance of the eagle? what the horses? what the varieties of coloured horses? what the rivers? what the chariots? what the wheels? what the so-called joy of the Angels?'' The printed ``anthromorphic'' is retained.}

The return to sensible figures completes the argument begun in chapter~II. Fire signifies a pervasive, illuminating, transforming activity; human features signify differentiated spiritual capacities. Eyes suggest knowledge, wings an unimpeded elevation, and teeth the division of unified intellectual nourishment into portions another can receive. Robes, instruments, winds, clouds, metals, animals, and wheels extend the same interpretive task. Their meanings arise from the relevant character and action of the image, not from an angelic anatomy assembled out of all the images at once.

Aquinas's use of the teeth is particularly instructive. In \ST{I q.~106 a.~1}, the image helps explain how a higher intellect adapts a more unified understanding to a lower recipient. The superior does not merely transmit a sentence: it makes truth intelligible according to the hearer's capacity. Elsewhere, the attribution of organs signifies powers without implying bodily sensation (\ST{I q.~51 a.~3 ad~2}). Chapter~XV also enters his distinction between illumination shared through the hierarchy and a communication addressed to a particular angel (\ST{I q.~107 a.~5, objection~3 and reply}). Imagery supplies questions that require further distinctions; it does not settle every question automatically.

Dionysius finally interprets angelic joy over human salvation without attributing bodily passion to the spirits. Their joy expresses a fitting participation in divine goodness and providential care. Even after this long catalogue he acknowledges what remains ``honoured by silence'' (XV.9, p.~66). That ending sets a limit to symbolic speculation. The work teaches an ordered reading of revealed signs and a return toward their divine source; it does not promise possession of every secret of heaven.
